\documentclass[10pt,a4paper]{amsart}
\usepackage[margin=19mm]{geometry}
\usepackage{amsmath,amssymb,mathtools}
\usepackage[scr=rsfs,cal=euler]{mathalfa}
\usepackage{xfrac}
\usepackage[unicode,hidelinks,bookmarks=false]{hyperref}
\newcommand{\ie}{i.e.\ }
\newcommand{\cf}{cf.\ }
\newcommand{\resp}{respectively\ }
\newcommand{\Subsection}{\subsection}
\providecommand{\nobreakdash}{\nobreak\hbox{-}}
\setlength{\emergencystretch}{2em}
\multlinegap0pt
\usepackage{bm}
\usepackage{bbm}
\usepackage{dsfont}
\usepackage{xspace}
\usepackage{cancel}
\usepackage{mathtools}
\usepackage{amsthm}
\makeatletter
\@ifundefined{theorem}{\newtheorem{theorem}{Theorem}}{}
\@ifundefined{proposition}{\newtheorem{proposition}[theorem]{Proposition}}{}
\makeatother
\title[Structured eigenbases and pair state transfer on threshold g]{Structured eigenbases and pair state transfer on threshold graphs}
\author{Leonardo de Lima, Renata Del-Vecchio, Hermie Monterde, Heber Teixeira}
\date{Source: arXiv:2601.13318v1 (CC BY 4.0)}
\begin{document}
\maketitle

\noindent\emph{Source and licence.} Structured eigenbases and pair state transfer on threshold graphs. Version arXiv:2601.13318v1 (CC BY 4.0), distributed under the Creative Commons Attribution 4.0 International licence, as recorded in the arXiv licence metadata for this exact version and shown on the abstract page https://arxiv.org/abs/2601.13318v1. Full rights notice: licenses/arxiv-2601.13318-CC-BY-4.0.txt.\par\smallskip

\noindent\textit{Editorial scope.} Counted unit: the Proposition whose source label is \texttt{prop:n-k} in the article (source0.tex lines 292--295, source label \texttt{prop:n-k}), delivered together with the article's own complete original proof of it (source0.tex lines 297--369, one \texttt{proof} environment of 73 lines). No mathematics is written, repaired or re-derived: statement and proof are the article's own text line for line. Required context retained uncounted: spectrum  threshold (lines 164--170); TeoPrincipal (lines 185--188). Every definition, display and label of those lines is retained unchanged. Omitted from this package and not redistributed: 1--163, 171--184, 189--291, 296--296, 370--1126. Nothing inside a retained excerpt is omitted or altered: every retained body is delivered exactly as the article's lines are written, so no omission receipt of any kind accompanies this package. Citations: the retained excerpts carry no citation command at all, so no bibliography entry of the article is transcribed and no reference is redistributed. Reference adaptation: none was needed and none was made; every reference inside the delivered unit resolves inside it, so no unresolved reference or citation is produced. Printed slips kept literal: no printed word, symbol, hypothesis or number of the counted statement or of its proof was altered. Layout and class: the article's own class is \texttt{elsarticle} with packages including hyperref, amsmath, amsthm, amssymb, color, comment, graphicx, enumitem, mathtools, multirow, multicol, lineno, bbm, xcolor; the wrapper is the lane's pinned \texttt{amsart} editorial wrapper at 10pt on A4, which restates the theorem-like environments the retained text needs and adds the article's own macro definitions that the retained text needs (none), with extra wrapper packages (bm, bbm, dsfont, xspace, cancel, mathtools). The delivered numbering is the unit's own. Assets: no retained span contains a figure environment, a graphics-inclusion command, a PDF-page inclusion, a table or a TikZ picture; no image file is copied into this package, the package declares no asset dependency and no image file is redistributed with it. Licence: version arXiv:2601.13318v1 of this article is distributed under the Creative Commons Attribution 4.0 International licence, as recorded in the arXiv licence metadata for this exact version and shown on the abstract page https://arxiv.org/abs/2601.13318v1. A full rights notice is delivered at licenses/arxiv-2601.13318-CC-BY-4.0.txt. Characters: every retained excerpt is pure ASCII, so the delivered engine, XeLaTeX with its default Latin Modern text font, reports no missing character.
    \begin{align}\label{spectrum  threshold}
    \mu_i = \left\{\begin{array}{rll}
        d_i +1,  & \text{if} \ \ 1 \leq i\leq \text{tr}(G),\\
        d_{i+1}, & \text{if} \ \ \text{tr}(G) + 1 \leq i \leq n-1,\\
        0,       & \text{if} \ \ i =n.
    \end{array}\right.
    \end{align}

\begin{theorem}
 \label{TeoPrincipal} 
Let $G$ be a connected graph on $n$ vertices. Then $G$ is a threshold graph if and only if  $\{\mathbf{x}^{1}, \ldots, \mathbf{x}^{n}\}$ is an eigenbasis for $L(G)$. 
\end{theorem}

\begin{proposition}\label{prop:n-k}
   Let $G$ be a connected threshold graph on $n$ vertices with binary sequence $\mathbf{b}=(b_1, \ldots, b_{l_2}, b_{l_2+1}, \ldots, b_{l_1-1}, b_{l_1}, b_{l_1+1}, \ldots, b_{n-1}, b_{n})$ such that $b_{l_1}=0$ and $b_{l_1+1}=\cdots=b_{n}=1$, for some $l_1 \in \{1,\ldots, n-1\}$. Let $l_2\in \{2, \ldots, l_1-1\}$ such that 
   $b_{l_2+1}=\cdots=b_{l_1}$ and $b_{l_2}\neq b_{l_2+1}$. Then $\mathcal{E}_{L}(n-l_1)$ is simply structured if and only if $2\leq l_2 \leq \lfloor \frac{l_1}{2} \rfloor$.
\end{proposition}

\begin{proof}
By Theorem \ref{TeoPrincipal}, $L(G)$ and $L(S_n)$ share the same eigenvectors. By hypothesis,  $b_{l_1}=0$ and $b_{l_1+1}=\cdots=b_{n}=1$, for some $l_1 \in \{1, \ldots, n-1\}$. Let $l_2 \in \{2, \ldots,l_1-1\}$ be such that $b_{l_2+1}=\cdots=b_{l_1}=0$ and $b_{l_2}=1$. Under these conditions, we have $\text{tr}(G)\geq n-l_1$ and from \eqref{spectrum  threshold}, it follows that $n-l_1$ is an eigenvalue of $L(G)$ with multiplicity $l_1-l_2$. By Theorem \ref{TeoPrincipal}, we have that  
\begin{align*}
    \mathcal{E}_L(n-l_1)&=\mathrm{span} \ \{ \mathbf{x}^{l_2}, \ldots, \mathbf{x}^{l_1-1} \} \quad  \textrm{and} \quad \mathrm{dim}(\mathcal{E}_{L}(n-l_1)) = l_1-l_2.
\end{align*} 
Suppose that $\mathcal{E}_L(n-l_1)$ is simply structured. Then, by hypothesis, there is a basis $B=\{\mathbf{u}^{l_2+1}, \ldots, \mathbf{u}^{l_1-1}, \mathbf{u}^{\prime}\}$ such that 
$$\mathcal{E}_L(n-l_1)=\mathrm{span} \ \{ \mathbf{x}^{l_2}, \ldots, \mathbf{x}^{l_1-1} \} = \mathrm{span}\{\mathbf{u}^{l_2+1}, \ldots, \mathbf{u}^{l_1-1},\mathbf{u}^{\prime}\},$$
where each vector in $B$ has entries only in the set $\{-1,0,1\}$.
Take $\mathbf{u}=(u_{1}, \ldots, u_{n}) \in B$ and determining it from the vectors $\mathbf{x}^{i}$ corresponds to finding coefficients $a=(a_{l_2}, a_{l_2+1}, \ldots, a_{l_1-1})^{T}$ such that 
\begin{equation}\label{eq:sistemalinear(n-k)}
    X a = \mathbf{u},
\end{equation}
where 
\begin{align*}
    X=[\mathbf{x}^{l_2}\cdots\mathbf{x}^{l_1-1}]=\left[\begin{array}{ccccc}
        1          & 1            & \ldots      & 1           & 1 \\ 
        \vdots     & \vdots       & \ddots      & \vdots      & \vdots \\ 
        1          & 1            & \ldots      & 1           & 1 \\
        -l_2         & 1            & \ldots      & 1           & 1 \\ 
        0          & -(l_2+1)       & \ldots      & 1           & 1 \\ 
        0          & 0            & \ldots      & 1           & 1 \\
        \vdots     & \vdots       & \ddots      & \vdots      & \vdots \\ 
        0          & 0            & \ldots      & 1           & 1 \\ 
        0          & 0            & \ldots      & -(l_1-2)    & 1 \\ 
        0          & 0            & \ldots      & 0           & -(l_1-1) \\
        0          & 0            & \ldots      & 0           & 0 \\ 
        \vdots     & \vdots       & \ddots      & \vdots      & \vdots \\ 
        0          & 0            & \ldots      & 0           & 0  
    \end{array}\right]. 
\end{align*}

Note that  $\mathbf{u}$ must have same number of $-1'$s and $1'$s since it is orthogonal to $\mathbbm{1}_n$. 
Observe that the first $l_2$ equations of \eqref{eq:sistemalinear(n-k)} are all identical, and therefore $u_{1}=u_{2}=\cdots=u_{l_2}.$ The other $l_1-l_2$ equations are given by 
\begin{equation}  \label{somas(n-k)}
\left\{\begin{array}{rrrrrrrrrrrr}
  -l_2 \, a_{l_2} + \hspace{1cm}  a_{l_2+1} + \hspace{0.2cm} \ldots \hspace{0.2cm} + \hspace{1cm} a_{l_1-2} + \hspace{1cm} a_{l_1-1} = & u_{l_2+1}   \\
                  -(l_2+1) \, a_{l_2+1} + \hspace{0.2cm} \ldots \hspace{0.2cm} + \hspace{1cm} a_{l_1-2} + \hspace{1cm} a_{l_1-1} = & u_{l_2+2}   \\
                                                                                       & \vdots   \\
                                         -(l_1-2) \, a_{l_1-2}  + \hspace{1cm} a_{l_1-1}   = & u_{l_1-1}   \\
                                                       -(l_1-1) \, a_{l_1-1}  = & u_{l_1}
\end{array}\right.
\end{equation}
and the remaining $n-l_1$ equations from  \eqref{eq:sistemalinear(n-k)} are all equal to 0, and therefore $u_{l_1+1}=\cdots=u_{n}=0$. 

Let us  analyze the cases where $u_1 = \cdots =u_{l_2} = 0$ or $u_1 = \cdots = u_{l_2} = 1.$ Note that the case where $u_1 = \cdots = u_{l_2} = -1$ is analogous to the previous case, except for the change of sign. 

First, suppose $u_1 = \cdots =u_{l_2} = 0.$ Since the number of $1$'s and $-1$'s are equal, for each $i \in \{l_2+1,\ldots,l_1-1\}$, take $u_{i} = 1, u_{i+1} = -1$ and $u_{j} = 0$ for all $j \ne \{i, i+1\}.$ The solution to the system \eqref{eq:sistemalinear(n-k)} is unique and given by 
\begin{align*}
    a_{i}=\frac{1}{i}, \ a_{i+1} = -\frac{1}{i} \ \text{and} \ a_{j}=0, \ \text{for} \  j \in \{l_2+1, \ldots, l_1-1\} \setminus \{i, i+1\}.    
\end{align*}
This shows that is possible to combine the vectors  $\mathbf{x}^{l_2+1},\ldots,\mathbf{x}^{l_1-1}$ to obtain $l_1-l_2-1$ eigenvectors of $L(G)$ that are linearly independent with entries in $\{-1,0,1\}$. These vectors are 
\begin{align}\label{ui estruturado(n)}
    \mathbf{u}^{i} = \frac{1}{i}\mathbf{x}^i-\frac{1}{i}\mathbf{x}^{i-1} = e_i-e_{i+1}
\end{align}
for each $i \in \{l_2+1,\ldots,l_1-1\}.$
Furthermore, in the other cases where are two or more $1'$s and $-1'$s, the vectors obtained are a linear combination of the vectors $\mathbf{u}^{i}$ obtained in \eqref{ui estruturado(n)}.

Now, suppose that $u_{1} = \cdots = u_{l_2} =1.$ In this case, we must have at least $l_2$ entries equal to $-1$ among $u_{l_2+1},\ldots,u_{l_1}$ in the $l_1-l_2$ equations in \eqref{somas(n-k)}. Hence, we have must $l_1 -l_2 \geq l_2,$ which implies that $2\leq l_2 \leq \lfloor \frac{l_1}{2} \rfloor.$ Suppose that $u_{l_2+1} = \cdots = u_{2l_2} =-1$ and let $\mathbf{u}^{\prime}=\sum_{i=1}^{l_2} e_i-\sum_{i=l_2+1}^{2l_2} e_{i}$. The solution to the system \eqref{eq:sistemalinear(n-k)} is unique and is given by 
\begin{align*}
    a_i=\frac{2l_2}{i(i+1)}, \ \text{for} \ \  i=l_2, \ldots, 2l_2-1 \ \text{and} \ a_i=0, \ \text{for} \ i=2l_2, \ldots, l_1-1.
\end{align*}
Thus, it is possible to combine the vectors $\mathbf{x}^{l_2},\ldots,\mathbf{x}^{2l_2-1}$ to obtain $1$ eigenvector of $L(G)$ with entries in $\{-1,0,1\}$ given by
\begin{align}\label{u' estruturado(n)}
    \mathbf{u}^{\prime}=\frac{2l_2}{l_2(l_2+1)} \, \mathbf{x}^{l_2} + \cdots + \frac{2l_2}{(2l_2-2)(2l_2-1)} \, \mathbf{x}^{2l_2-2} + \frac{1}{2l_2-1} \, \mathbf{x}^{2l_2-1}
\end{align}
Note that $\mathbf{u}^{\prime}$ in \eqref{u' estruturado(n)} is not a linear combination of the vectors  $\mathbf{u}^{l_2+1}, \ldots, \mathbf{u}^{l_1-1}$ in \eqref{ui estruturado(n)}, since the first $l_2$ coordinates of $\mathbf{u}^{\prime}$ are equal to 1, while the first $l_2$ coordinates of $\mathbf{u}^{l_2+1}, \ldots, \mathbf{u}^{l_1-1}$ are equal to 0. Therefore, the vectors in the set $\{\mathbf{u}^{l_2+1},\ldots,\mathbf{u}^{l_1-1}, \mathbf{u}^{\prime}\}$ are linearly independent. Thus, if $\mathcal{E}_L(n-l_1)$ is simply structured, then $2\leq l_2 \leq \lfloor \frac{l_1}{2} \rfloor.$

On the other hand, suppose that $2\leq l_2 \leq \lfloor \frac{l_1}{2} \rfloor$ and that $b_{l_2+1}=\cdots=b_{l_1}=0$ and $b_{l_2}=1$. So,
\begin{center}
$\mathcal{E}_L(n-l_1)=\mathrm{span} \ \{ \mathbf{x}^{l_2}, \ldots, \mathbf{x}^{l_1-1} \} \quad  \textrm{and} \quad \left\lceil \frac{l_1}{2} \right\rceil \leq \mathrm{dim}(\mathcal{E}_{L}(n-l_1)) \leq l_1-2.$
\end{center} 
Since $l_2 \in \{2,\ldots, \lfloor \frac{l_1}{2} \rfloor\}$, it is possible to combine the vectors $\mathbf{x}^{l_2}, \ldots, \mathbf{x}^{l_1-1}$ that are in the basis of $\mathcal{E}_L(n-l_1)$ to obtain a new basis for $\mathcal{E}_L(n-l_1)$ formed by the vectors $\mathbf{u}^{l_2+1}, \ldots, \mathbf{u}^{l_1-1},\mathbf{u}^{\prime}$ with entries only in $\{-1,0,1\}$, where the vectors $\mathbf{u}^{i}$, for $i=l_2+1, \ldots, l_1-1$ are given in \eqref{ui estruturado(n)} and the vector $\mathbf{u}^{\prime}$ is given in \eqref{u' estruturado(n)}. Therefore, if $2\leq l_2 \leq \lfloor \frac{l_1}{2} \rfloor$, then $\mathcal{E}_L(n-l_1)$ is simply structured.
\end{proof} 

\end{document}
